\itemhead{习题4.1 · 第4题（教材p.104）}{XJ-B1-C04-HW-P104-Y0874}
\label{XJ-B1-C04-HW-P104-Y0874}
{\footnotesize 湘教版教材p.104 / PDF第111页 / 习题4.1 / 第4题}
\field{原题}{计算（式中字母都表示正数）：
\subq{1}$x^{1/3}\cdot x^{3/4}\cdot x^{-1/12}$；
\subq{2}$a^{3/2}\cdot a^{2/3}\cdot a^{-5/6}$；
\subq{3}$\left(x^{1/2}y^{-1/3}\right)^6$；
\subq{4}$\left(\dfrac{4s^2r^{-6}}{9t^4}\right)^{3/2}$；
\subq{5}$6x^{2/3}y^{-1/3}\div\left(\dfrac32x^{-4/3}y^{2/3}\right)$；
\subq{6}$\left(x^2-2+x^{-2}\right)\div\left(x^2-x^{-2}\right)$；
\subq{7}$\dfrac{x-y}{x^{2/3}+x^{1/3}y^{1/3}+y^{2/3}}$；
\subq{8}$\dfrac{a^{1/2}-b^{1/2}}{a^{1/2}+b^{1/2}}-\dfrac{a^{1/2}+b^{1/2}}{a^{1/2}-b^{1/2}}$。}
\field{参考答案}{依次为 $x,a^{4/3},x^3/y^2,8s^3/(27r^9t^6),4x^2/y,(x^2-1)/(x^2+1),x^{1/3}-y^{1/3},-4\sqrt{ab}/(a-b)$。第（6）仍须 $x\ne1$；第（8）须 $a\ne b$，均沿用原式限制。}
{\small \field{关联知识}{P02 有理数、实数指数幂与运算（教材帮p.144、145 / PDF153、154）；F02 定义域与原式限制（教材帮p.79、80、81 / PDF88、89、90）}
\field{方法与易错点}{统一为指数幂，分开处理系数与各底数的指数；最终复查定义条件及数值精度。 易错：负指数误作负数；幂的幂指数相加；约分丢禁点；过早四舍五入。}
\field{答案依据}{答案册相关题组：PDF第31页（册内页30）；按当前题库原题逐项重算，不按旧版题号机械对套。}}
